\section{Related Work}

\paragraph{Quantum programming languages and substructural types.}
Quantum programming languages use linear and affine type disciplines to reflect no-cloning and controlled disposal of quantum data.
Recent POPL work on Qurts uses Rust-like lifetimes to control automatic uncomputation through a substructural discipline~\cite{hirata2025qurts}.
\system shares the view that quantum programming needs static ownership information, but its ownership discipline applies to encoded code states and fault-tolerant resources rather than to automatic uncomputation.
In particular, a $\res{TState}$ in \system is treated as a linear physical artifact whose consumption is part of the target acquisition plan.

\paragraph{Resource estimation for quantum circuit descriptions.}
Type-based resource estimation for quantum circuit description languages uses effects, refinements, and index terms to derive circuit-size bounds~\cite{colledan2025resource}.
\system also uses effects, but the quantities are specialized to fault-tolerant acquisition: logical error, failure, acceptance, code switches, factories, decoder latency, physical qubits, cycles, and qubit-rounds.
The effect system is therefore tied to rule certificates and backend side conditions, not only to circuit size.
This distinction is visible in the case study, where two plans can implement the same logical gate sequence but expose different switch/factory/certificate profiles.

\paragraph{Quantum compiler generation and mapping.}
Compiler frameworks for qubit mapping and routing identify common state-machine structure across diverse architectures and synthesize compilers from architecture specifications~\cite{molavi2026qmr}.
\system is complementary.
It does not attempt to solve low-level mapping and routing optimally; instead, it exposes code capabilities and fault-tolerant acquisition decisions as typed program structure.
The current geometry checker is a conservative rectangular embedding and workspace-capacity check, not a full nearest-neighbor routing or lattice-surgery scheduler.

\paragraph{Parameterized quantum verification.}
Automata-based verification can prove relational properties of parameterized quantum circuits across all input sizes~\cite{abdulla2026parameterized}.
\system targets a different verification layer.
Rather than verifying arbitrary circuit families against functional specifications, it checks that a compiler-generated fault-tolerant acquisition plan respects code capabilities, linear resources, certificate levels, backend constraints, and conservative effect bounds.
The two approaches are compatible: parameterized circuit verification can supply facts about circuit families, while \system can record how such facts are used in an acquisition plan.

\paragraph{Fault-tolerant protocols.}
\system relies on the standard magic-state view of universal fault-tolerant computation~\cite{bravyi2005magic} and recent protocol results for code switching, magic-state distillation, and surface-code implementation profiles~\cite{butt2024switching,heussen2025oneway,itogawa2025zero,wan2024constant}.
The calculus does not replace those protocol proofs.
It provides a language-level composition layer that records which protocol facts are locally checked, which are imported, and which are still assumptions in the prototype rule library.
